\documentclass[11pt]{article}

\usepackage{amsmath,amssymb,amsthm,mathtools,bm}
\usepackage{booktabs}
\usepackage{geometry}
\usepackage{enumitem}
\usepackage{float}
\usepackage{hyperref}
\usepackage{xcolor}
\usepackage{graphicx}

\geometry{margin=1in}

% ------------------------------------------------------------
% Draft status
% ------------------------------------------------------------
% Updated after the latest algorithmic development (2026-08-26).
% The main algorithmic line in this draft is now:
%   admissible convexity-preserving kinetic-term regularization
%   + vacuum-compatible potential regularization
%   + positivity/mass-preserving Fourier pseudospectral discretization
%   + FISTA-CD with potential-proximal splitting
%   + interior Newton stationarity refinement.
%
% Numerical tables remain placeholders unless a value has already been
% certified by the current code.  No unrun convergence claim is made.
% ------------------------------------------------------------

\newcommand{\R}{\mathbb{R}}
\newcommand{\C}{\mathbb{C}}
\newcommand{\A}{\mathcal{A}}
\newcommand{\E}{\mathcal{E}}
\newcommand{\CN}{\mathcal{C}_N}
\newcommand{\SN}{S_N}
\newcommand{\SNP}{S_N^+}
\newcommand{\dd}{\,\mathrm{d}}
\newcommand{\ii}{\mathrm{i}}
\newcommand{\eps}{\varepsilon}
\newcommand{\sig}{\sigma}
\newcommand{\inner}[2]{\left\langle #1,#2\right\rangle}
\newcommand{\norm}[1]{\left\|#1\right\|}
\newcommand{\abs}[1]{\left|#1\right|}
\newcommand{\grad}{\nabla}
\newcommand{\half}{\frac12}
\newcommand{\todo}[1]{\textcolor{red}{[TODO: #1]}}

\theoremstyle{plain}
\newtheorem{theorem}{Theorem}[section]
\newtheorem{lemma}[theorem]{Lemma}
\newtheorem{proposition}[theorem]{Proposition}
\newtheorem{corollary}[theorem]{Corollary}

\theoremstyle{definition}
\newtheorem{definition}[theorem]{Definition}
\newtheorem{remark}[theorem]{Remark}

\numberwithin{equation}{section}
\counterwithin{figure}{section}
\counterwithin{table}{section}

\title{\bf
A positive-conservative Fourier optimization method for ground states of Bose--Einstein condensates with higher-order interactions
}

\author{
Bo Lin,\thanks{{\color{red} Institute}} \quad
Xinran Ruan\thanks{(Corresponding author) School of Mathematical Sciences, Capital Normal University, Beijing 100048, People's Republic of China {\tt (xinran.ruan@cnu.edu.cn)}.}
}
\date{}

\begin{document}
\maketitle

\begin{abstract}
We develop a positive-conservative Fourier optimization (PCFO) method
for computing ground states of Bose--Einstein condensates with
higher-order interactions. The ground-state problem admits a convex
density formulation, but the singular behavior near zero density
poses difficulties for high-accuracy computation. We introduce
convex regularizations of the density energy and establish their
$\Gamma$-convergence on the computational domain as the regularization
parameters vanish. The regularized problem is discretized by a
Fourier pseudospectral method with mass conservation and nodal
positivity. 
For the resulting constrained optimization problem, we use an accelerated
first-order method for the main energy reduction, followed by a Newton
refinement to reduce the remaining optimality residual. 
Numerical experiments show spectral-type convergence
for regularized problems, quantify the effects of the
regularization parameters on accuracy and spatial resolution, and
demonstrate the effectiveness of the two-stage optimization method.
\end{abstract}

\noindent
\textbf{Key words.}
Bose--Einstein condensate; higher-order interaction; density
formulation; convex optimization; Fourier
pseudospectral method; positivity preserving. 

\medskip

\noindent
\textbf{2020 MSC.}
35Q55, 65N35, 65K10, 90C25.


% ============================================================
\section{Introduction}
% ============================================================

Bose--Einstein condensation, first realized experimentally in dilute
atomic gases in 1995 \cite{Anderson1995}, provides a fundamental
setting for the study of macroscopic quantum phenomena. 
In the dilute and low-temperature regime, Bose--Einstein condensates are commonly
described by the Gross--Pitaevskii equation \cite{BaoCaiReview}, in
which two-body interactions are modeled by a zero-range contact
potential. Finite-range and effective-range corrections arise beyond
this approximation
\cite{CollinMassignanPethick2007,FuWangGao2003} and may become more
relevant under strong confinement or near narrow Feshbach resonances
\cite{Chin2010,VekslerFishmanKetterle2014,ZinnerThogersen2009}.
One such correction leads to a modified Gross--Pitaevskii model with
a higher-order interaction involving the gradient of the density.

We consider the modified Gross--Pitaevskii model with higher-order
interactions (HOI), whose energy is
\begin{equation}
\label{eq:wave-energy}
E(\phi)
=
\int_{\R^d}
\left[
\half |\grad\phi|^2
+V(x)|\phi|^2
+\frac{\beta}{2}|\phi|^4
+\frac{\delta}{2}\bigl|\grad(|\phi|^2)\bigr|^2
\right]\dd x,
\end{equation}
subject to
\begin{equation}
\label{eq:wave-mass}
\norm{\phi}_{L^2}^2=1.
\end{equation}
Here $V(x)\ge0$ is a confining potential, $\beta\ge 0$ is the repulsive
contact-interaction strength, and $\delta>0$ is the higher-order
interaction strength. The competition between the contact and
higher-order interactions can affect ground-state structures and
their asymptotic regimes \cite{BaoCaiRuan2019,RuanCaiBao2016}. Under
strong confinement, the higher-order interaction also modifies the
effective lower-dimensional mean-field models
\cite{RuanCaiBao2016}.
%In this work, we focus on the
%repulsive regime $\beta\ge0$.

As shown in \cite{BaoCaiRuan2019}, a ground state may be chosen real
and nonnegative. Introducing the density $\rho=|\phi|^2$ gives the
equivalent density formulation
\begin{equation}
\label{eq:density-energy}
\E_0(\rho)
=
\int_{\R^d}
\left[
\frac{|\grad\rho|^2}{8\rho}
+V(x)\rho
+\frac{\beta}{2}\rho^2
+\frac{\delta}{2}|\grad\rho|^2
\right]\dd x,
\end{equation}
over
\begin{equation}
\label{eq:admissible-density}
\A
=
\left\{
\rho\in H^1(\R^d):
\rho\ge0,\quad
\int_{\R^d}\rho\,\dd x=1,\quad
\int_{\R^d}V\rho\,\dd x<\infty
\right\}.
\end{equation}
The quotient $|\grad\rho|^2/(8\rho)$ is interpreted in the extended
sense defined in Section~\ref{sec:density}. For $\beta\ge0$,
$\E_0$ is strictly convex on $\A$. This convexity is useful for
ground-state computation, although the kinetic term remains singular
at zero density.


Ground-state computation for the Gross--Pitaevskii equation has been
studied extensively. Early approaches include direct energy
minimization and normalized gradient-flow methods
\cite{BaoDu2004,BaoTang2003}, together with spectrally accurate
discretizations \cite{BaoChernLim2006}. For strongly nonlinear or
rotating condensates, Sobolev-gradient methods, preconditioned
spectral solvers, nonlinear conjugate-gradient methods, regularized
Newton methods, and Riemannian optimization have been developed
\cite{AntoineDuboscq2014,AntoineLevittTang2017,
DanailaKazemi2010,DanailaProtas2017,WuWenBao2017}. More recent work
includes normalized Sobolev-gradient flows and nonlinear eigenvalue
iterations with convergence analysis
\cite{AltmannHenningPeterseim2021,ChenLuLuZhang2024,
HenningPeterseim2020}. Fourier pseudospectral discretizations combined
with manifold optimization have also been applied to more general
spinor condensates \cite{TianCaiWuWen2020}.

For the HOI model, a normalized gradient-flow method with an
attractive--repulsive splitting was developed in \cite{Ruan2018}.
W. Bao and X. Ruan \cite{BaoRuan2019} subsequently formulated the
ground-state problem as a convex minimization problem in the density
variable, using a shifted kinetic regularization, a second-order
finite-difference discretization, and an accelerated projected-gradient
method. A square-root regularization of the potential was also
considered there as a numerical enhancement.

The present work builds on this convex density formulation with the
aim of improving spatial accuracy and final stationarity. We extend
the kinetic regularization to a class of convexity-preserving
denominators and treat the potential regularization with an
independent scale. The regularized problem is discretized by a Fourier
pseudospectral method with exact mass conservation and nodal
positivity. For the resulting constrained optimization problem, we
use a potential-proximal accelerated first-order method followed by a
Newton refinement. We establish the $\Gamma$-convergence of the
regularized energies on the computational domain and examine the
spatial accuracy, regularization effects, and solver performance
numerically.

The rest of the paper is organized as follows.
Section~\ref{sec:density} introduces the convex regularizations and
establishes their variational convergence.
Section~\ref{sec:ps} presents the Fourier pseudospectral discretization
and the positive mass-conservative projection.
Section~\ref{sec:opt} develops the potential-proximal FISTA method and
the Newton refinement. Section~\ref{sec:num} reports the numerical
results, and Section~\ref{sec:summary} concludes the paper.

% ============================================================
\section{Density formulation and convex regularization}
\label{sec:density}
% ============================================================

In this section, we study the convex structure of the density energy,
introduce convex regularizations of the kinetic and potential terms,
and establish variational convergence of the regularized energies.


\subsection{Convex regularization of the kinetic and potential terms}

\paragraph{Kinetic regularization.}

We extend the density kinetic term to possible vacuum regions.
For $s\ge0$ and $q\in\R^d$, let
\begin{equation}
\label{eq:kinetic-integrand}
F_0(s,q)
=
\begin{cases}
\dfrac{|q|^2}{8s}, & s>0,\\[1ex]
0, & s=0,\ q=0,\\
+\infty, & s=0,\ q\ne0.
\end{cases}
\end{equation}
With this convention, the density energy in
\eqref{eq:density-energy} can be written as
\begin{equation}
\label{eq:density-energy-extended}
\E_0(\rho)
=
\int_{\R^d}
\left[
F_0(\rho,\nabla\rho)
+V\rho
+\frac{\beta}{2}\rho^2
+\frac{\delta}{2}|\nabla\rho|^2
\right]\dd x,
\qquad
\rho\in\A.
\end{equation}

The density energy has the following convexity property.

\begin{proposition}[Convexity of the density energy]
\label{prop:density-convexity}
Assume $\beta\ge0$ and $\delta>0$.
Then $\E_0$ is strictly convex on $\A$ and hence has at most one
minimizer.
\end{proposition}

\begin{proof}
The map $(s,q)\mapsto |q|^2/s$ is the perspective of the convex
quadratic function $q\mapsto |q|^2$, and
\eqref{eq:kinetic-integrand} is its lower-semicontinuous convex
extension to $s=0$. Hence the kinetic term is convex. The potential
term is linear, and the $\beta$ term is convex for $\beta\ge0$.

The $\delta$ term gives strict convexity on $\A$. Indeed, equality in
its convexity inequality for two densities in $\A$ implies that their
gradients coincide almost everywhere. Their difference is therefore
constant, and integrability on $\R^d$ implies that this constant is
zero.
\end{proof}

To remove the singularity at zero density while preserving the convex
structure, we regularize the kinetic denominator. For $s\ge0$ and
$q\in\R^d$, set
\begin{equation}
\label{eq:regularized-kinetic-integrand}
F_\eps(s,q)
=
\frac{|q|^2}{8r_\eps(s)},
\end{equation}
where $r_\eps$ is a positive regularization of the denominator.

\begin{definition}[Admissible denominator]
\label{def:admissible-regularization}
A family $\{r_\eps\}_{\eps>0}$ is called admissible if, for every
$\eps>0$,
\begin{equation}
\label{eq:regularization-admissibility-local}
r_\eps\in C^1([0,\infty)),
\qquad
r_\eps>0,
\qquad
r_\eps'\ge0,
\qquad
r_\eps \text{ is concave},
\end{equation}
and, for every $s\ge0$,
\begin{equation}
\label{eq:regularization-admissibility-limit}
s
\le
r_{\eps_1}(s)
\le
r_{\eps_2}(s)
\quad
\text{for }0<\eps_1\le\eps_2,
\qquad
r_\eps(s)\to s
\quad\text{as }\eps\to0^+.
\end{equation}
\end{definition}

The concavity of the denominator preserves the
convexity of the regularized kinetic term.

\begin{lemma}[Convexity of the regularized kinetic integrand]
\label{lem:regularized-kinetic-convexity}
Let $r_\eps$ be admissible in the sense of Definition \ref{def:admissible-regularization}. Then
\[
F_\eps(s,q)
=
\frac{|q|^2}{8r_\eps(s)}
\]
is jointly convex in $(s,q)\in[0,\infty)\times\R^d$.
\end{lemma}

\begin{proof}
Let $0\le\theta\le1$ and set
\[
s_\theta=\theta s_1+(1-\theta)s_2,
\qquad
q_\theta=\theta q_1+(1-\theta)q_2.
\]
By the concavity of $r_\eps$,
\[
r_\eps(s_\theta)
\ge
\theta r_\eps(s_1)+(1-\theta)r_\eps(s_2).
\]
Hence
\[
F_\eps(s_\theta,q_\theta)
\le
\frac{|q_\theta|^2}
{8[\theta r_\eps(s_1)+(1-\theta)r_\eps(s_2)]}.
\]
The weighted Cauchy--Schwarz inequality gives
\[
\frac{|q_\theta|^2}
{\theta r_\eps(s_1)+(1-\theta)r_\eps(s_2)}
\le
\theta\frac{|q_1|^2}{r_\eps(s_1)}
+
(1-\theta)\frac{|q_2|^2}{r_\eps(s_2)},
\]
which proves the result.
\end{proof} 

The simplest example is the shifted denominator introduced in
\cite{BaoRuan2019},
\begin{equation}
\label{eq:shift-denominator}
r_\eps(s)=s+\eps.
\end{equation}
For comparison, we also consider the family
\begin{equation}
\label{eq:piecewise-regularization-family}
r_{\eps,m}(s)
=
\begin{cases}
s+\eps-\dfrac{\eps}{m+1}
\left(1-\dfrac{s}{\eps}\right)^{m+1},
&0\le s\le\eps,\\[2ex]
s+\eps,
&s>\eps,
\end{cases}
\qquad m=1,2.
\end{equation}
These functions belong to $C^m([0,\infty))$ and satisfy
Definition~\ref{def:admissible-regularization}. The cases $m=1,2$
will be included in the numerical comparison later in Section \ref{sec:num}.

\paragraph{Potential regularization.}

The kinetic regularization removes the singularity at $\rho=0$, but
the derivative of the linear potential term $V\rho$ at zero density
is $V$. In the discrete problem, this can activate the positivity
constraint and reduce the regularity relevant to the Fourier
approximation. We therefore regularize the potential term near zero
density by replacing $\rho$ with
\begin{equation}
\label{eq:potential-regularizer}
p_\sig(s)
=
\sqrt{s^2+\sig^2}-\sig,
\qquad s\ge0.
\end{equation}
A square-root smoothing of this type was considered in
\cite{BaoRuan2019} as a numerical enhancement. Here $\sig$ is treated
as an independent regularization scale and is analyzed together with
the kinetic regularization. We have
\begin{equation}
\label{eq:potential-regularizer-derivatives}
p_\sig(0)=p_\sig'(0)=0,
\qquad
p_\sig''(s)
=
\frac{\sig^2}{(s^2+\sig^2)^{3/2}}
\ge0.
\end{equation}
Thus $Vp_\sig$ is convex for $V\ge0$. The condition
$p_\sig'(0)=0$ removes the nonzero derivative of the potential term
at zero density, while
$p_\sig''(0)=1/\sig$ explains the increasing curvature as $\sig$
decreases and motivates the proximal treatment in
Section~\ref{subsec:potential_FISTA}.

Moreover, $p_\sig(s)\to s$ as $\sig\to0^+$. For the convergence
analysis below, it is sufficient to take
$\sig=\sig(\eps)\to0$ as $\eps\to0^+$.

Combining the two regularizations, we obtain
\begin{equation}
\label{eq:joint-regularized-energy}
\E_{\eps,\sig}(\rho)
=
\int_{\R^d}
\left[
F_\eps(\rho,\nabla\rho)
+V(x)p_\sig(\rho)
+\frac{\beta}{2}\rho^2
+\frac{\delta}{2}|\nabla\rho|^2
\right]\dd x,
\qquad
\rho\in\A.
\end{equation}

The convex structure of the original density problem is retained by
the regularized energy.

\begin{proposition}[Convexity of the regularized energy]
\label{prop:joint-convexity}
Assume that $r_\eps$ is admissible in the sense of Definition \ref{def:admissible-regularization}, $V\ge0$, $\beta\ge0$,
$\delta>0$, and $\sig>0$.
Then $\E_{\eps,\sig}$ \eqref{eq:joint-regularized-energy} is strictly convex on $\A$ \eqref{eq:admissible-density}.
If $\beta>0$, it is strongly convex with respect to the $L^2$ norm.
\end{proposition}

\begin{proof}
By Lemma~\ref{lem:regularized-kinetic-convexity}, the kinetic
contribution is convex. The potential term is convex because
$p_\sig$ is convex and $V\ge0$, while the remaining terms are convex.
Strict convexity follows from the $\delta$ term as in
Proposition~\ref{prop:density-convexity}. If $\beta>0$, the term
$\frac{\beta}{2}\|\rho\|_{L^2}^2$ gives strong convexity in $L^2$.
\end{proof}




\subsection{\texorpdfstring{$\Gamma$}{Gamma}-convergence of the regularized energies}

Let
\[
D_L=[-L,L)^d
\]
be the bounded periodic computational domain used for the Fourier
discretization. We study the $\Gamma$-convergence of the regularized
energies on $D_L$ as the regularization parameters vanish, with
respect to the strong $L^2(D_L)$ topology
\cite{Braides2002}.

Define
\begin{equation}
\label{eq:admissible-density-box}
\A_L
=
\left\{
\rho\in H^1_{\rm per}(D_L):
\rho\ge0\ {\rm a.e.},
\quad
\int_{D_L}\rho\,\dd x=1
\right\}.
\end{equation}
For $\rho\in\A_L$, define
\begin{equation}
\label{eq:energies-DL}
\begin{aligned}
\E_{0,L}(\rho)
&=
\int_{D_L}
\left[
F_0(\rho,\nabla\rho)
+V\rho
+\frac{\beta}{2}\rho^2
+\frac{\delta}{2}|\nabla\rho|^2
\right]\dd x,\\
\E_{\eps,\sig,L}(\rho)
&=
\int_{D_L}
\left[
F_\eps(\rho,\nabla\rho)
+Vp_\sig(\rho)
+\frac{\beta}{2}\rho^2
+\frac{\delta}{2}|\nabla\rho|^2
\right]\dd x.
\end{aligned}
\end{equation}

The potential regularization is uniformly small on $D_L$. Since
\begin{equation}
\label{eq:potential-energy-defect}
0\le s-p_\sig(s)\le\sig,
\qquad s\ge0,
\end{equation}
we have
\[
\left|
\int_{D_L}
V\,[p_\sig(\rho)-\rho]\,\dd x
\right|
\le
\sig\int_{D_L}V\,\dd x.
\]
Thus the potential contribution is a uniformly vanishing perturbation
as $\sig\to0$, and the main convergence argument can be carried out
for the kinetic regularization.

\begin{theorem}[$\Gamma$-convergence of the joint regularization]
\label{thm:gamma-joint}
Let $V\in C(D_L)$ with $V\ge0$, $\beta\ge0$, and $\delta>0$.
Let $\{r_\eps\}$ be admissible in the sense of
Definition~\ref{def:admissible-regularization}, and let
$\sig(\eps)>0$ satisfy
$\sig(\eps)\to0$ as $\eps\to0^+$.
Extend $\E_{\eps,\sig(\eps),L}$ and $\E_{0,L}$ to $L^2(D_L)$
by setting them equal to $+\infty$ outside $\A_L$.
Then
\[
\E_{\eps,\sig(\eps),L}
\xrightarrow{\Gamma}
\E_{0,L}
\qquad
\text{in the strong }L^2(D_L)\text{ topology}.
\]
Moreover, the family
$\{\E_{\eps,\sig(\eps),L}\}_{\eps>0}$
is equicoercive in $L^2(D_L)$.
\end{theorem}

\begin{proof}
We first consider the kinetic regularization alone and keep the
potential term unchanged. Define
\[
\widetilde\E_{\eps,L}(\rho)
=
\int_{D_L}
\left[
F_\eps(\rho,\nabla\rho)
+V\rho
+\frac{\beta}{2}\rho^2
+\frac{\delta}{2}|\nabla\rho|^2
\right]\dd x.
\]

We first prove the liminf inequality. Let $\eps_n\to0$ and
$\rho_n\to\rho$ strongly in $L^2(D_L)$. If
\[
\liminf_{n\to\infty}
\widetilde\E_{\eps_n,L}(\rho_n)=+\infty,
\]
there is nothing to prove. Otherwise, after passing to a subsequence,
we may assume that the energies are uniformly bounded. Since
$\delta>0$, the gradient term together with the fixed mass gives a
uniform $H^1(D_L)$ bound. Hence, after extracting a further
subsequence,
\[
\rho_n\rightharpoonup\rho
\qquad
\text{weakly in }H^1(D_L).
\]

Fix $\bar\eps>0$. For all sufficiently large $n$,
$\eps_n\le\bar\eps$. The ordering in
\eqref{eq:regularization-admissibility-limit} therefore gives
\[
F_{\eps_n}(s,q)
\ge
F_{\bar\eps}(s,q).
\]
For fixed $\bar\eps$, the integrand $F_{\bar\eps}$ is continuous and
jointly convex in $(s,q)$. The associated integral functional is
therefore weakly lower semicontinuous on $H^1(D_L)$. Together with
the remaining terms, this yields
\[
\liminf_{n\to\infty}
\widetilde\E_{\eps_n,L}(\rho_n)
\ge
\widetilde\E_{\bar\eps,L}(\rho).
\]
Finally, let $\bar\eps\downarrow0$. Since
$r_{\bar\eps}(s)\downarrow s$, the nonnegative kinetic integrands
increase pointwise to $F_0$. The monotone convergence theorem then
gives
\[
\liminf_{n\to\infty}
\widetilde\E_{\eps_n,L}(\rho_n)
\ge
\E_{0,L}(\rho).
\]

For the recovery sequence, let $\rho\in\A_L$ satisfy
$\E_{0,L}(\rho)<\infty$ and take the constant sequence
$\rho_\eps=\rho$. The same monotone convergence argument gives
\[
\widetilde\E_{\eps,L}(\rho)
\longrightarrow
\E_{0,L}(\rho)
\qquad
\text{as }\eps\to0^+.
\]
Hence
\[
\widetilde\E_{\eps,L}
\xrightarrow{\Gamma}
\E_{0,L}
\qquad
\text{in strong }L^2(D_L).
\]

It remains to restore the potential regularization.
By \eqref{eq:potential-energy-defect},
\[
\sup_{\rho\in\A_L}
\left|
\E_{\eps,\sig(\eps),L}(\rho)
-
\widetilde\E_{\eps,L}(\rho)
\right|
\le
\sig(\eps)\int_{D_L}V\,\dd x
\longrightarrow0.
\]
Thus the potential smoothing is a uniformly vanishing perturbation
and does not change the $\Gamma$-limit. 

The equicoercivity follows from the $\delta$ term and the fixed mass,
which give a uniform $H^1(D_L)$ bound, together with the compact
embedding $H^1(D_L)\hookrightarrow L^2(D_L)$.
\end{proof}

The $\Gamma$-convergence result gives the corresponding convergence
of minimum values and minimizers.

\begin{corollary}[Convergence of minimizers]
\label{cor:min-conv}
Under the assumptions of Theorem~\ref{thm:gamma-joint}, let
$\rho_{\eps,L}$ and $\rho_{0,L}$ denote the minimizers of
$\E_{\eps,\sig(\eps),L}$ and $\E_{0,L}$ on $\A_L$, respectively.
Then
\begin{equation}
\label{eq:minimizer-convergence}
\min_{\rho\in\A_L}\E_{\eps,\sig(\eps),L}(\rho)
\longrightarrow
\min_{\rho\in\A_L}\E_{0,L}(\rho),
\qquad
\rho_{\eps,L}
\longrightarrow
\rho_{0,L}
\quad\text{strongly in }L^2(D_L).
\end{equation}
\end{corollary}

\begin{proof}
We first note that the minimizers are well defined and unique.
Existence follows from the standard direct method of the calculus of
variations. Strict convexity follows from the $\delta$ term together
with the mass constraint. Indeed, equality in the convexity
inequality implies that two densities in $\A_L$ differ by a constant,
and their equal masses imply that this constant is zero.

With uniqueness established, Theorem~\ref{thm:gamma-joint} and
equicoercivity give the convergence of the minimum values and the
subsequential convergence of $\rho_{\eps,L}$. Every convergent
subsequence has a minimizer of $\E_{0,L}$ as its limit. Since this
minimizer is unique, all such limits coincide with $\rho_{0,L}$, and
therefore the whole family converges.
\end{proof}

% ============================================================
\section{Fourier pseudospectral discretization}
\label{sec:ps}
% ============================================================

We discretize the regularized density problem on a periodic
computational domain by a Fourier pseudospectral method. We introduce
the discrete energy and its derivatives, together with the positive
mass-conservative projection used in the constrained optimization.

We present the one-dimensional formulation on $D_L=[-L,L)$; the
multidimensional case follows by the standard tensor-product
extension. Let $N$ be even,
$x_j=-L+jh$, $j=0,\ldots,N-1$, and $h=2L/N$. For
$u,v\in\R^N$, define the discrete inner product and norm by
\begin{equation}
\label{eq:h-inner}
\inner{u}{v}_h
=
h\sum_{j=0}^{N-1}u_jv_j,
\qquad
\norm{u}_h^2=\inner{u}{u}_h.
\end{equation}


\subsection{Discrete energy, gradient, and Hessian structure}

Let $D_N$ denote the real Fourier pseudospectral first-derivative
operator. For even $N$, we set the derivative of the Fourier mode
$k=N/2$ to zero, so that $D_N$ maps real grid vectors to real grid
vectors and satisfies
\begin{equation}
\label{eq:skew-adjoint}
D_N^T=-D_N
\end{equation}
with respect to the inner product \eqref{eq:h-inner}. The action of
$D_N$ is evaluated by FFTs.

For a grid density $\rho$, let $q=D_N\rho$. The discrete regularized
energy is
\begin{equation}
\label{eq:discrete-energy}
E_{\eps,\sig,N}(\rho)
=
h\sum_{j=0}^{N-1}
\left[
\frac{q_j^2}{8r_\eps(\rho_j)}
+V_jp_\sig(\rho_j)
+\frac{\beta}{2}\rho_j^2
+\frac{\delta}{2}q_j^2
\right],
\end{equation}
subject to the discrete positivity and mass constraints
\begin{equation}
\label{eq:discrete-feasible}
\CN
=
\left\{
\rho\in\R^N:
\rho_j\ge0,
\quad
h\sum_{j=0}^{N-1}\rho_j=1
\right\}.
\end{equation}
Proposition~\ref{prop:discrete-convex} shows that the discrete energy
\eqref{eq:discrete-energy} retains the convexity of the regularized
problem.

\begin{proposition}[Discrete convexity]
\label{prop:discrete-convex}
Let $r_\eps$ be admissible in the sense of
Definition~\ref{def:admissible-regularization}. Let $\sig>0$ and assume
$V_j\ge0$, $\beta\ge0$, and $\delta>0$.
Then $E_{\eps,\sig,N}$ is convex on $\CN$ and admits a minimizer.
If $\beta>0$, it is strongly convex with respect to
$\norm{\cdot}_h$, and the minimizer is unique.
\end{proposition}


\begin{proof}
By Lemma 2.3 and the linearity of
\[
\rho\mapsto \bigl(\rho_j,(D_N\rho)_j\bigr),
\]
the discrete kinetic term is convex. The remaining terms inherit
the convexity properties established in Proposition 2.4.
Since $\CN$ is compact and $E_{\eps,\sig,N}$ is continuous, a
minimizer exists. If $\beta>0$, the term
$\frac{\beta}{2}\|\rho\|_h^2$ gives strong convexity and hence
uniqueness.
\end{proof}

Differentiating \eqref{eq:discrete-energy} gives the discrete gradient.
Let $q=D_N\rho$. With all nonlinear operations understood
pointwise, the gradient with respect to the inner product
\eqref{eq:h-inner} is characterized by
\begin{equation}
\label{eq:discrete-gradient}
\begin{aligned}
&G_{\eps,\sig,N}(\rho)
=
\frac14D_N^T
\left(
\frac{q}{r_\eps(\rho)}
\right)
-\frac18
\frac{q^2r_\eps'(\rho)}
     {r_\eps(\rho)^2}
+V\,p_\sig'(\rho)
+\beta\rho
+\delta D_N^TD_N\rho,\\
&\left.
\frac{\dd}{\dd t}
E_{\eps,\sig,N}(\rho+t\eta)
\right|_{t=0}
=
\inner{G_{\eps,\sig,N}(\rho)}{\eta}_h .
\end{aligned}
\end{equation}

At points where $r_\eps$ is twice differentiable, define
\begin{equation}
\label{eq:kinetic-hessian-factorization}
\begin{aligned}
B
&=
D_N-
\operatorname{diag}
\left(
\frac{q\,r_\eps'(\rho)}
     {r_\eps(\rho)}
\right),
&
W
&=
\operatorname{diag}
\left(
\frac{1}{4r_\eps(\rho)}
\right).
\end{aligned}
\end{equation}
The Hessian contribution of the regularized kinetic term admits the
factorization
\begin{equation}
\label{eq:kinetic-hessian}
H_{\rm kin}
=
B^TWB
+
\operatorname{diag}
\left(
-\frac{q^2r_\eps''(\rho)}
       {8r_\eps(\rho)^2}
\right).
\end{equation}
Since $r_\eps>0$ and $r_\eps''\le0$ wherever the second derivative
exists, both terms in \eqref{eq:kinetic-hessian} are positive
semidefinite. The Hessian of the full discrete energy therefore has
the decomposition
\begin{equation}
\label{eq:full-hessian-spd-form}
H
=
H_{\rm kin}
+\delta D_N^TD_N
+\beta I
+\operatorname{diag}\!\bigl(Vp_\sig''(\rho)\bigr).
\end{equation}
Hence $H$ is symmetric positive semidefinite. If $\beta>0$, it is
positive definite. This structure will be used in the Newton refinement in
Section~\ref{subsec:Newton-PCG}.

For the shifted denominator $r_\eps(s)=s+\eps$, which is used in the
Newton refinement below, $r_\eps'=1$ and $r_\eps''=0$. The diagonal
correction in \eqref{eq:kinetic-hessian} then vanishes. The
corresponding matrix-free Hessian action is recorded in
Appendix~\ref{app:newton-implementation}.

\subsection{Positive mass-conservative projection}

The constrained optimization requires a projection that preserves
nonnegativity and mass. Let $\SN$ denote the trigonometric
interpolation space associated with the Fourier grid, and define
\begin{equation}
\label{eq:SNplus}
\SNP
=
\left\{
v_N\in\SN:
v_N(x_j)\ge0,\quad j=0,\ldots,N-1
\right\}.
\end{equation}
where, for $v_N\in\SN$, we use the same discrete norm notation
$\norm{v_N}_h^2
=
h\sum_{j=0}^{N-1}|v_N(x_j)|^2$.
For $z_N\in\SN$, consider the positive mass-conservative
pseudospectral projection \cite{lin2024sinum}
\begin{equation}
\label{eq:spectral-positive-proj}
\operatorname*{argmin}_{g_N\in\SNP}
\frac12\norm{g_N-z_N}_h^2
\quad
\text{subject to}
\quad
\int_{D_L}g_N\,\dd x
=
\int_{D_L}z_N\,\dd x .
\end{equation}
The following proposition gives its nodal form.

\begin{proposition}[Nodal form of the positive mass-conservative projection]
\label{prop:proj-equiv}
Let $z=(z_j)_{j=0}^{N-1}\in\R^N$, where $z_j=z_N(x_j)$.
Then \eqref{eq:spectral-positive-proj} is equivalent to
\[
\operatorname*{argmin}_{y\in\R^N}
\frac12\norm{y-z}_h^2
\]
subject to
\[
y_j\ge0,
\qquad
h\sum_{j=0}^{N-1}y_j
=
\int_{D_L}z_N\,\dd x .
\]
\end{proposition}

\begin{proof}
By the definition of the discrete norm,
\[
\norm{g_N-z_N}_h^2
=
h\sum_{j=0}^{N-1}
|g_N(x_j)-z_N(x_j)|^2.
\]
Moreover, the trapezoidal rule is exact for trigonometric
polynomials, so
\[
\int_{D_L}g_N\,\dd x
=
h\sum_{j=0}^{N-1}g_N(x_j).
\]
Hence the objective and constraints in
\eqref{eq:spectral-positive-proj} coincide with those of the nodal
problem.
\end{proof}

Proposition~\ref{prop:proj-equiv} shows that the pseudospectral
projection can be carried out entirely in terms of the nodal values.
For the unit-mass constraint used in the optimization, we therefore
define the nodal projection onto $\CN$ by
\begin{equation}
\label{eq:nodal-projection}
P_{\CN}(z)
=
\operatorname*{argmin}_{y\in\CN}
\frac12\norm{y-z}_h^2,
\qquad z\in\R^N.
\end{equation}
It is given componentwise by
\begin{equation}
\label{eq:simplex-form}
\bigl[P_{\CN}(z)\bigr]_j
=
\max\{z_j-\lambda,0\},
\qquad
h\sum_{j=0}^{N-1}
\bigl[P_{\CN}(z)\bigr]_j
=
1,
\end{equation}
where $\lambda$ is uniquely determined by the mass constraint.

% ============================================================
\section{Two-stage optimization method}
\label{sec:opt}
% ============================================================

In this section, we solve the discrete convex minimization problem
using a two-stage strategy. We first apply an accelerated proximal
method to reduce the energy efficiently, and then refine the resulting
state through the KKT system to reduce the stationarity residual.

The discrete ground-state problem is
\begin{equation}
\label{eq:discrete-min}
\min_{\rho\in\CN}E_{\eps,\sig,N}(\rho).
\end{equation}
Since the objective and $\CN$ are convex, every KKT point is a global
minimizer.

To assess stationarity in both stages, for a fixed $\tau_r>0$ define
the projected-gradient residual
\begin{equation}
\label{eq:pg-map}
\mathcal G_{\tau_r}(\rho)
=
\frac{1}{\tau_r}
\left[
\rho
-
P_{\CN}
\left(
\rho-\tau_rG_{\eps,\sig,N}(\rho)
\right)
\right].
\end{equation}
For any feasible $\rho$, the optimality condition for
\eqref{eq:discrete-min} is equivalent to
$\mathcal G_{\tau_r}(\rho)=0$. We use
$\norm{\mathcal G_{\tau_r}(\rho)}_h$ as the stationarity measure.


\subsection{Potential-proximal FISTA}
\label{subsec:potential_FISTA}

FISTA is an accelerated proximal-gradient method for composite convex
optimization \cite{BeckTeboulle2009}. We apply it to the discrete
problem \eqref{eq:discrete-min} by splitting the energy into a smooth part and a proximal
part.
As shown in  \eqref{eq:potential-regularizer-derivatives}, the potential regularization has curvature $p_\sig''(0)=1/\sig$, which
becomes large as $\sig$ decreases. 
We therefore include the regularized
potential together with the positivity and mass constraints in the
proximal part. Thus we use the splitting
\begin{equation}
\label{eq:f-plus-g}
E_{\eps,\sig,N}(\rho)
=
f_{\eps,N}(\rho)+g_{\sig,N}(\rho),
\end{equation}
where
\begin{equation}
\label{eq:smooth-part}
f_{\eps,N}(\rho)
=
h\sum_j
\left[
\frac{(D_N\rho)_j^2}{8r_\eps(\rho_j)}
+\frac{\beta}{2}\rho_j^2
+\frac{\delta}{2}(D_N\rho)_j^2
\right]
\end{equation}
is the smooth part, and
\begin{equation}
\label{eq:prox-part}
g_{\sig,N}(\rho)
=
h\sum_j V_jp_\sig(\rho_j)
+
I_{\CN}(\rho)
\end{equation}
is the proximal part. Here $I_{\CN}$ denotes the indicator function
of $\CN$.
The gradient of $f_{\eps,N}$ with respect to the discrete inner
product \eqref{eq:h-inner} is
\begin{equation}
\label{eq:smooth-gradient}
\nabla_h f_{\eps,N}(\rho)
=
\frac14D_N^T
\left(
\frac{D_N\rho}{r_\eps(\rho)}
\right)
-\frac18
\frac{(D_N\rho)^2r_\eps'(\rho)}
     {r_\eps(\rho)^2}
+\beta\rho
+\delta D_N^TD_N\rho,
\end{equation}
where the nonlinear operations are understood pointwise.

We use the Chambolle--Dossal variant of FISTA with backtracking
\cite{BeckTeboulle2009, ChambolleDossal2015}.
Given two consecutive iterates $\rho^{k-1},\rho^k\in\CN$, one FISTA
iteration consists of an extrapolation to $y^k$, a forward step to
$z^k$, and a proximal update to $\rho^{k+1}\in\CN$.

The extrapolated state is
\begin{equation}
\label{eq:fista-extrapolation}
y^k
=
\rho^k+\theta_k(\rho^k-\rho^{k-1}),
\end{equation}
where $\theta_k$ is the Chambolle--Dossal extrapolation parameter.
The extrapolation preserves the mass constraint. Since
$f_{\eps,N}$ is defined for nonnegative densities, the extrapolation
is accepted only if $y^k\ge0$. Otherwise, it is restarted by setting
\[
y^k=\rho^k.
\]

For a trial step size $\tau_k>0$, the forward step is
\begin{equation}
\label{eq:fista-forward}
z^k
=
y^k-\tau_k\nabla_h f_{\eps,N}(y^k).
\end{equation}
The proximal update is
\begin{equation}
\label{eq:potential-prox}
\rho^{k+1}
=
\operatorname*{argmin}_{\rho\in\CN}
\left\{
\frac12\norm{\rho-z^k}_h^2
+
\tau_k h\sum_jV_jp_\sig(\rho_j)
\right\}.
\end{equation}
The positivity and mass constraints are therefore enforced in the
proximal step.

The trial step size is accepted when the standard majorization
condition
\begin{equation}
\label{eq:fista-backtracking}
\begin{aligned}
f_{\eps,N}(\rho^{k+1})
\le\;&
f_{\eps,N}(y^k)
+
\inner{\nabla_h f_{\eps,N}(y^k)}
       {\rho^{k+1}-y^k}_h 
+
\frac{1}{2\tau_k}
\norm{\rho^{k+1}-y^k}_h^2
\end{aligned}
\end{equation}
is satisfied. Otherwise, $\tau_k$ is reduced and the
forward--proximal step is repeated.

The proximal problem \eqref{eq:potential-prox} can be solved through
a single scalar mass multiplier. For a fixed $\lambda$, each positive
component of $\rho^{k+1}$ satisfies
\begin{equation}
\label{eq:potential-prox-node}
s-z_j^k
+\tau_kV_jp_\sig'(s)
+\lambda
=
0,
\qquad s>0.
\end{equation}
The left-hand side is strictly increasing in $s$, since
\(
1+\tau_kV_jp_\sig''(s)>0.
\)
Hence the positive solution, when it exists, is unique. The $j$th
component is zero when $z_j^k\le\lambda$ and otherwise is given by
the unique positive solution of
\eqref{eq:potential-prox-node}. The multiplier $\lambda$ is determined
by the mass constraint
\[
h\sum_j\rho_j^{k+1}=1.
\]
Thus the proximal update consists of independent scalar nodal
equations coupled only through the mass constraint.

For the unregularized linear potential, $p(s)=s$,
\eqref{eq:potential-prox} reduces to
\[
\rho^{k+1}
=
P_{\CN}(z^k-\tau_kV).
\]
Hence the potential-proximal update is a nonlinear extension of the
nodal projection \eqref{eq:nodal-projection}.

Once the energy has stabilized and the projected residual \eqref{eq:pg-map} is
sufficiently small, the iterate is passed to the Newton refinement.


\subsection{Newton refinement}
\label{subsec:Newton-PCG}

After the FISTA stage, we use a Newton refinement to further reduce
the stationarity residual \eqref{eq:pg-map}. The corresponding KKT conditions of
\eqref{eq:discrete-min} are 
\begin{equation}
\label{eq:kkt}
\begin{aligned}
& G_{\eps,\sig,N}(\rho)+\lambda\bm1-\mu=0,
\qquad
h\bm1^T\rho=1,\\
& \rho\ge0,\qquad
\mu\ge0,\qquad
\rho_j\mu_j=0.
\end{aligned}
\end{equation}
When the solution is interior, the inequality multiplier vanishes and
\eqref{eq:kkt} reduces to
\begin{equation}
\label{eq:interior-kkt}
G_{\eps,\sig,N}(\rho)+\lambda\bm1=0,
\qquad
h\bm1^T\rho=1.
\end{equation}

Starting from the FISTA output, each interior Newton iteration
computes a correction from the linearized KKT system and applies a
damped update. In the strongly convex regime $\beta>0$, the Hessian
is symmetric positive definite by
\eqref{eq:full-hessian-spd-form}. For the current iterate
$(\rho^k,\lambda^k)$, define
\[
r^k
=
G_{\eps,\sig,N}(\rho^k)+\lambda^k\bm1,
\qquad
r_m^k
=
h\bm1^T\rho^k-1
\]
as the stationarity and mass residuals, respectively.
Let $H^k=H(\rho^k)$ denote the Hessian of
$E_{\eps,\sig,N}$ at $\rho^k$ \eqref{eq:full-hessian-spd-form}.
The Newton correction
$(d^k,d_\lambda^k)$ satisfies
\[
H^kd^k+d_\lambda^k\bm1=-r^k,
\qquad
h\bm1^Td^k=-r_m^k.
\]

Rather than solving the indefinite KKT system directly, we use a
Schur complement. Let
\[
H^kz^k=r^k,
\qquad
H^kw^k=\bm1.
\]
Then
\begin{equation}
\label{eq:schur-two-solves}
d_\lambda^k
=
\frac{r_m^k-h\bm1^Tz^k}
     {h\bm1^Tw^k},
\qquad
d^k
=
-z^k-d_\lambda^kw^k.
\end{equation}
Thus each interior Newton step requires two linear systems with the
same symmetric positive definite Hessian. We solve these systems by
the preconditioned conjugate-gradient (PCG) method.

The correction is applied with a damping factor
$\alpha_k\in(0,1]$:
\[
\rho^{k+1}
=
\rho^k+\alpha_kd^k,
\qquad
\lambda^{k+1}
=
\lambda^k+\alpha_kd_\lambda^k.
\]
The damping parameter is chosen to preserve positivity and to
globalize the Newton iteration. If active constraints are
encountered, a standard primal--dual active-set treatment can be used 
\cite{Ulbrich2011}.

For the PCG solves, we use a sparse variable-coefficient
finite-difference approximation of the Fourier pseudospectral
Hessian as a symmetric positive definite preconditioner. Its
construction is given in
Appendix~\ref{app:newton-implementation}.


\section{Numerical experiments}
\label{sec:num}

In this section, we examine the spatial accuracy, regularization
effects, and computational performance of the positive-conservative
Fourier optimization (PCFO) method, and present representative
two-dimensional examples.

Unless otherwise stated, we take $\beta=\delta=10$, impose the
unit-mass constraint, use the shifted denominator
$r_\eps(\rho)=\rho+\eps$, and adopt the potential regularization
\eqref{eq:potential-regularizer}.
% ------------------------------------------------------------
\subsection{Spatial accuracy}
\label{subsec:num-mesh}
% ------------------------------------------------------------

We first examine the spatial accuracy of the Fourier discretization
with fixed regularization parameters. 
To separate the spatial error from
the regularization error, we work on the periodic interval
$D_L=[-L,L)$ with $L=32$ and fix
$\eps=\sig=10^{-2}$. 

We take the harmonic potential $V_H(x)=x^2/2$. Its periodic extension
is continuous but not smoothly periodic, since
$V_H'(-L)\ne V_H'(L)$, which limits the Fourier accuracy at high
resolution. For the spatial-convergence test, we therefore use
\begin{equation}
\label{eq:periodic-harmonic}
V_L^{\rm per}(x)
=
\bigl(1-\chi_L(x)\bigr)\frac{x^2}{2}
+
\chi_L(x)\frac{L^2}{2},
\end{equation}
where $\chi_L$ is an even $C^\infty$ cutoff satisfying
$\chi_L=0$ for $|x|\le0.75L$ and $\chi_L=1$ for $|x|\ge0.90L$.
Thus $V_L^{\rm per}$ agrees with the harmonic potential in the
interior and is constant near the periodic boundary.

Let $\rho_N$ and $\rho_{\rm ref}$ denote the trigonometric
interpolants of the numerical solutions on the $N$-point grid and
the finest comparison grid, respectively. Let $P_N$ denote the
$L^2(D_L)$-orthogonal Fourier projection onto $\SN$. We define
\begin{equation}
\label{eq:spatial-errors}
\begin{aligned}
e_{\rho,2}(N)
&=
\norm{\rho_N-\rho_{\rm ref}}_{L^2(D_L)},\quad
e_{\rho,\infty}(N)
&=
\norm{\rho_N-P_N\rho_{\rm ref}}_{\ell^\infty}.
\end{aligned}
\end{equation}
The energy error is
\begin{equation}
\label{eq:spatial-energy-error}
e_E(N)
=
\abs{
E_{\eps,\sig,N}(\rho_N)
-
E_{\eps,\sig,N_{\rm ref}}(\rho_{\rm ref})
}.
\end{equation}



\begin{table}[H]
\centering
\small
\caption{
Fourier mesh refinement for $\eps=\sig=10^{-2}$ using the harmonic
trap with a $C^\infty$ periodic continuation \eqref{eq:periodic-harmonic}.
The comparison state is computed with $N_{\rm ref}=8192$.
}
\label{tab:mesh}
\begin{tabular}{ccccccc}
\toprule
$N$
&
$e_{\rho,2}$
&
rate
&
$e_{\rho,\infty}$
&
rate
&
$e_E$
&
rate
\\
\midrule
32
& $2.519\times10^{-2}$
& --
& $1.282\times10^{-2}$
& --
& $5.831\times10^{-3}$
& --
\\
64
& $3.394\times10^{-3}$
& $2.89$
& $1.422\times10^{-3}$
& $3.17$
& $2.077\times10^{-3}$
& $1.49$
\\
128
& $1.190\times10^{-4}$
& $4.83$
& $9.022\times10^{-5}$
& $3.98$
& $5.597\times10^{-6}$
& $8.54$
\\
256
& $6.227\times10^{-7}$
& $7.58$
& $2.526\times10^{-7}$
& $8.48$
& $3.770\times10^{-9}$
& $10.54$
\\
512
& $6.415\times10^{-11}$
& $13.24$
& $4.185\times10^{-11}$
& $12.56$
& $4.174\times10^{-14}$
& --
\\
\bottomrule
\end{tabular}
\end{table}

Table~\ref{tab:mesh} reports the errors relative to the
$N_{\rm ref}=8192$ comparison state. The state and energy errors
decrease rapidly, with increasing empirical orders, until the
numerical floor is reached, which is consistent with spectral-type
spatial convergence for fixed regularization parameters.

For comparison, we repeat the refinement with the unmodified harmonic
potential $V_H=x^2/2$ on the periodic domain. The two calculations are
nearly indistinguishable on coarse and moderate meshes. 
At higher resolutions, the direct harmonic calculation develops an
algebraic error floor, with the observed order approaching two,
consistent with the lack of periodic smoothness of $V_H$.
This error
floor is removed by the $C^\infty$ far-field continuation, as shown
in Fig.~\ref{fig:mesh-convergence}.

\begin{figure}[H]
\centering
\includegraphics[width=0.7\textwidth]
{figs/mesh_convergence_1d_spatial_error.eps}
\caption{
Fourier spatial convergence with fixed regularization parameters.
``Continued harmonic'' refers to the $C^\infty$ far-field
continuation \eqref{eq:periodic-harmonic}, and ``direct harmonic'' to
$V_H(x)=x^2/2$ on the periodic computational domain.
}
\label{fig:mesh-convergence}
\end{figure}


% ------------------------------------------------------------
\subsection{Regularization effects}
\label{subsec:num-regularization}
% ------------------------------------------------------------

The regularized model contains two parameters associated with the
kinetic and potential terms. We first examine the effect of the
potential regularization parameter $\sig$ at fixed $\eps$, including
its influence on Fourier resolution. We then compare different
regularizations of the density kinetic term as $\eps\to0$.

\paragraph{Potential regularization.}

We fix $\eps=10^{-2}$, $L=32$, and $N=4096$, and take
$\sig=\eps^p$ with $p=1,2,3,4$. Let $\rho_{\eps,\sig}$ denote the
minimizer with the regularized potential, and let $\rho_{\eps,0}$
denote the minimizer obtained with 
the original potential. To quantify the perturbation introduced
by $p_\sig$, we define
\begin{equation}
\label{eq:num-potential-bias}
\begin{aligned}
B_E(\sig)
&=
E_{\eps,0}(\rho_{\eps,\sig})
-
E_{\eps,0}(\rho_{\eps,0}),
\quad
B_\rho(\sig)
&=
\norm{
\rho_{\eps,\sig}-\rho_{\eps,0}
}_{L^2(D_L)}.
\end{aligned}
\end{equation}
Here $E_{\eps,0}$ denotes the energy with the same denominator
$r_\eps$ and the original potential contribution $V\rho$.



Table~\ref{tab:sigma-sweep} shows that both perturbations decrease
approximately linearly with $\sig$ over the tested range. The minimum
density exhibits the same scaling for small $\sig$, indicating that
the positive low-density background introduced by the potential
regularization vanishes together with $\sig$.

\begin{table}[H]
\centering
\small
\caption{
Effect of the potential regularization scale at
$\eps=10^{-2}$, $L=32$, and $N=4096$.
The quantities $B_E(\sig)$ and $B_\rho(\sig)$ are defined in
\eqref{eq:num-potential-bias}.
}
\label{tab:sigma-sweep}
\begin{tabular}{cccc}
\toprule
$\sig$
&
$B_E(\sig)$
&
$B_\rho(\sig)$
&
$\min\rho$
\\
\midrule
$10^{-2}$
& $1.776\times10^{0}$
& $2.548\times10^{-2}$
& $6.449\times10^{-5}$
\\
$10^{-4}$
& $1.849\times10^{-2}$
& $2.900\times10^{-4}$
& $6.740\times10^{-7}$
\\
$10^{-6}$
& $1.849\times10^{-4}$
& $2.719\times10^{-6}$
& $6.743\times10^{-9}$
\\
$10^{-8}$
& $1.849\times10^{-6}$
& $1.993\times10^{-8}$
& $6.744\times10^{-11}$
\\
\bottomrule
\end{tabular}
\end{table}

To examine the transition introduced by the potential regularization,
we consider
\begin{equation}
\label{def:p_prime}
p_\sig'\bigl(\rho_{\eps,\sig}(x)\bigr)
=
\frac{\rho_{\eps,\sig}(x)}
{\sqrt{\rho_{\eps,\sig}(x)^2+\sig^2}},
\end{equation}
which measures the local slope of the regularized potential relative
to the original linear potential. It is close to one when
$\rho_{\eps,\sig}\gg\sig$ and close to zero when
$\rho_{\eps,\sig}\ll\sig$.

Figure~\ref{fig:regularization-sigma} shows that the transition
between these two regimes becomes sharper as $\sig$ decreases.
Accordingly, the modification introduced by the potential
regularization becomes more localized near the low-density region.
The resulting smaller spatial scale requires more Fourier modes to
resolve accurately.

%The profiles show that the potential regularization mainly modifies
%the low-density region while leaving the bulk region essentially
%unchanged. This behavior is consistent with the condition
%$p_\sig'(0)=0$, which removes this nonzero derivative at zero density
%while localizing the modification to the low-density region. As
%$\sig$ decreases, this modification
%becomes increasingly localized near the transition region. The
%resulting narrower transition introduces a smaller spatial scale,
%which requires more Fourier modes to resolve accurately.

\begin{figure}[t]
\centering
\includegraphics[width=0.70\textwidth]
{figs/regularization_sigma_L32_transition.eps}
\caption{
Profiles of
$p_\sig'(\rho_{\eps,\sig}(x))$ \eqref{def:p_prime}
for different potential-regularization scales.
The transition becomes sharper as $\sig$ decreases, and the inset
enlarges the transition region.
All computations use
$\eps=10^{-2}$, $L=32$, and $N=4096$.
}
\label{fig:regularization-sigma}
\end{figure}



To examine this effect on Fourier resolution, we repeat the mesh
refinement for
$\sig=0$, $10^{-2}$, $10^{-4}$, and $10^{-6}$, while keeping
$\eps=10^{-2}$ and $L=32$. Here $\sig=0$ corresponds to the original
linear potential without potential regularization. We use the same
$C^\infty$-continued harmonic potential
\eqref{eq:periodic-harmonic} and the $N_{\rm ref}=65536$ solution as
the reference.

Figure~\ref{fig:sigma-accuracy} shows the resulting energy and state
errors. For $\sig=10^{-2}$, both errors rapidly enter the
spectral-type regime and reach their numerical floors. As $\sig$
decreases, a longer pre-asymptotic range appears before the rapid
decay becomes visible. In contrast, the case $\sig=0$ does not exhibit
spectral-type decay over the tested resolutions, and the errors
decrease much more slowly. 
This reveals a bias--resolution tradeoff in which decreasing $\sig$
reduces the regularization error but requires a finer Fourier mesh to
maintain the same spatial accuracy.


\begin{figure}[H]
\centering
\includegraphics[width=0.48\textwidth]
{figs/sigma_effect_energy_error_1d.eps}
\includegraphics[width=0.48\textwidth]
{figs/sigma_effect_state_error_1d.eps}
\caption{
Effect of the potential regularization scale on Fourier accuracy.
The left panel shows the energy error and the right panel shows the
total spectral $L^2$ state error. Each curve refines a fixed-$\sig$
problem with $\eps=10^{-2}$, $L=32$, and the harmonic-core
$C^\infty$ far-field continuation, relative to the
$N_{\rm ref}=65536$ comparison state.
}
\label{fig:sigma-accuracy}
\end{figure}




\paragraph{Regularization of the kinetic term.}

We next show the effect of the kinetic regularization by fixing the
potential regularization and the spatial discretization. We compare
the shifted denominator \eqref{eq:shift-denominator} with the
piecewise $C^1$ and $C^2$ denominators in
\eqref{eq:piecewise-regularization-family}. Throughout this test, we
take $\sig=10^{-2}$, $L=32$, and $N=512$, and use the same
$C^\infty$-continued harmonic potential. The tested values are
$\eps=10^{-1},5\times10^{-2},2\times10^{-2},10^{-2},
5\times10^{-3},2\times10^{-3},10^{-3}$.

A common reference state $\rho_{\rm ref}$ is computed with the shifted
denominator at $\eps_{\rm ref}=10^{-5}$. To compare the three
regularizations on the same basis, we measure the state error and
evaluate the energy error with the common unregularized kinetic
denominator $r_0(\rho)=\rho$. 
We define
\begin{equation}
\label{eq:kinetic-reg-errors}
\begin{aligned}
e_\rho(\eps)
&=
\norm{\rho_\eps-\rho_{\rm ref}}_{L^2(D_L)},
\quad
e_E(\eps)
&=
\abs{
E_{0,\sig}(\rho_\eps)
-
E_{0,\sig}(\rho_{\rm ref})
}.
\end{aligned}
\end{equation}
Here $E_{0,\sig}$ denotes the energy with the original kinetic
denominator $r_0(s)=s$ and the fixed potential regularization
$p_\sig$. 

\begin{figure}[H]
\centering
\includegraphics[width=0.48\textwidth]
{figs/kinetic_regularization_energy_error.eps}
\hfill
\includegraphics[width=0.48\textwidth]
{figs/kinetic_regularization_state_error.eps}
\caption{
Effect of the kinetic regularization. The shifted denominator \eqref{eq:shift-denominator} and the
piecewise $C^1$ and $C^2$ denominators \eqref{eq:piecewise-regularization-family} are compared using the errors
defined in \eqref{eq:kinetic-reg-errors}. The left panel shows the
energy error \(e_\rho(\eps)\) and the right panel the $L^2$ state error \(e_E(\eps)\). The dashed and
dotted lines indicate first- and second-order slopes, respectively.
All computations use $\sig=10^{-2}$, $L=32$, and $N=512$.
}
\label{fig:kinetic-regularization}
\end{figure}

Figure~\ref{fig:kinetic-regularization} shows that the three
regularizations have essentially the same asymptotic behavior as
$\eps$ decreases. 
The common-functional energy error approaches second order whereas the state error is approximately first order.
The piecewise $C^1$ denominator gives a modest reduction in error for
the larger values of $\eps$, but neither piecewise regularization
shows a clear asymptotic rate advantage over the shifted denominator.
This supports using the simpler shifted denominator in the solver
experiments.
%In all cases, the final projected-gradient residual is well below the
%reported regularization errors.

% ------------------------------------------------------------
\subsection{Solver performance}
\label{subsec:num-solver}
% ------------------------------------------------------------

We next examine the computational performance of the two-stage
optimization strategy. Throughout this subsection, we use the shifted
denominator $r_\eps(\rho)=\rho+\eps$. We first examine the
respective roles of FISTA and the Newton refinement, and then
examine how the two regularization scales affect the computational
cost.

\paragraph{FISTA and Newton refinement.}

We consider the representative case
$\eps=10^{-2}$, $\sig=10^{-4}$, $L=32$, and $N=4096$.
All projected-gradient residuals \eqref{eq:pg-map} in this subsection use $\tau_r=1$.
The FISTA stage is switched to the Newton refinement when
\(
\norm{\mathcal G_1(\rho^k)}_h\le10^{-5}.
\)
The criterion is checked after the first $200$ FISTA iterations and
is required to hold for two consecutive iterations. Starting from the
same switching state, one branch continues with FISTA, while the other
proceeds with the Newton refinement. The two potential-proximal runs
therefore share the same FISTA history up to the switch.

For comparison, we also consider a smooth-potential FISTA, in which
the regularized potential is included in the smooth part instead of the proximal part. 
It minimizes the same
discrete energy and uses the same FISTA backtracking rule.

Since FISTA and Newton steps have different computational costs,
Fig.~\ref{fig:solver-performance} reports the convergence histories
against cumulative wall-clock time. 
The two branches share the same FISTA time history up to the switching
point. Thereafter, the additional time of each branch is measured from
the common switching time.
The left panel shows the energy error
relative to the final two-stage solution, while the right panel shows
the projected-gradient residual. 
As shown in Fig.~\ref{fig:solver-performance}, FISTA rapidly reduces
the energy and produces an accurate approximation of the ground state,
but continued FISTA decreases the projected-gradient residual only
slowly. Starting from the same switching state, the Newton refinement
reduces the residual much more rapidly.

\begin{figure}[H]
\centering
\includegraphics[width=0.48\textwidth]
{figs/solver_energy_history_shared_prefix.eps}
\hfill
\includegraphics[width=0.48\textwidth]
{figs/solver_residual_history_shared_prefix.eps}
\caption{
Convergence histories for continued FISTA and the two-stage
FISTA--Newton method with
$\eps=10^{-2}$, $\sig=10^{-4}$, $L=32$, and $N=4096$.
The left panel shows the energy error relative to the final two-stage
solution, and the right panel shows the projected-gradient residual
$\norm{\mathcal G_1(\rho^k)}_h$.
Here ``potential-proximal FISTA'' and ``smooth-potential FISTA''
refer to placing the regularized potential in the proximal and smooth
parts of the splitting, respectively.
}
\label{fig:solver-performance}
\end{figure}




\paragraph{Dependence on the regularization scales.}

We next examine how the computational cost changes as the
regularization parameters decrease. To isolate the two effects, we
first vary $\eps$ at fixed $\sig=10^{-2}$ and then vary $\sig$ at
fixed $\eps=10^{-2}$. In all cases,
$r_\eps(\rho)=\rho+\eps$.

All runs use the same Fourier resolution $N=4096$, initial density,
and solver parameters. Thus the comparison isolates the effect of the
regularization scales at a fixed Fourier resolution. The results are summarized in
Table~\ref{tab:solver-regularization-cost}.
The reported wall-clock time includes both the FISTA and Newton
stages.

\begin{table}[H]
\centering
\caption{
Computational cost of the two-stage solver for different
regularization scales. The shifted denominator
$r_\eps(\rho)=\rho+\eps$ is used in all cases, with the same
initialization, Fourier resolution $N=4096$, and solver parameters.
``FISTA'' and ``Newton'' denote the iteration counts in the two
stages, and ``PCG'' is the accumulated number of inner PCG
iterations. The last column reports the final projected-gradient
residual $\|\mathcal G_1(\rho)\|_h$ \eqref{eq:pg-map}. The asterisk
indicates that the FISTA stage reached the prescribed iteration limit.
}
\label{tab:solver-regularization-cost}
\begin{tabular}{ccccccc}
\toprule
$\eps$ & $\sig$
& FISTA & Newton & PCG
& time (s) & $\|\mathcal G_1(\rho)\|_h$ \\
\midrule
\input{tables/solver_regularization_cost_rows.tex}\\
\bottomrule
\end{tabular}
\end{table}

The FISTA stage is limited to $20000$ iterations. At this limit,
the Newton refinement is entered only if
$\|\mathcal G_1(\rho)\|_h\le10^{-4}$. The asterisk in
Table~\ref{tab:solver-regularization-cost} marks the run that reached
this limit.

At fixed $\sig=10^{-2}$, decreasing $\eps$ increases the coefficients
involving $1/r_\eps(\rho)$ in the kinetic Hessian, particularly in
low-density regions. The FISTA iteration counts and total solution
times increase over the tested values in
Table~\ref{tab:solver-regularization-cost}. The smallest $\eps$ case
reaches the prescribed FISTA iteration limit and has a somewhat larger
final residual than the other runs.

By contrast, decreasing $\sig$ at fixed $\eps=10^{-2}$ produces only
a modest increase in the total cost over the tested range. Although a
smaller $\sig$ sharpens the potential transition and changes the cost
of the proximal and Newton stages, the
potential-proximal splitting keeps the large curvature
$p_\sig''(0)=1/\sig$ outside the smooth FISTA gradient. This behavior
is consistent with the motivation for the splitting in
Subsection~\ref{subsec:potential_FISTA}.


% ------------------------------------------------------------
\subsection{Two-dimensional examples}
\label{subsec:num-2d}
% ------------------------------------------------------------

We present two examples in two dimensions to illustrate the
applicability of the method to more general trapping potentials.

Throughout this subsection, we use the shifted denominator
$r_\eps(\rho)=\rho+\eps$, the potential regularization $p_\sig$, and
the parameters
$\beta=\delta=10$, $\eps=10^{-2}$, and $\sig=10^{-4}$.
The computations are carried out on the periodic box
$[-16,16]^2$ with $N_x=N_y=512$ Fourier collocation points.
As in the one-dimensional experiments, the harmonic part of the
potential is left unchanged in the central region and smoothly
continued to a constant near the periodic boundary.

We first consider the anisotropic harmonic trap
\begin{equation}
\label{eq:2d-anisotropic-harmonic}
V_{\rm ah}(x,y)
=
\frac12
\left(
\gamma_x^2x^2+\gamma_y^2y^2
\right),
\qquad
\gamma_x=1,
\quad
\gamma_y=2.
\end{equation}
The resulting ground-state density is shown in the left panel of
Fig.~\ref{fig:two-dimensional-examples}. The stronger confinement in
the $y$-direction produces an anisotropic density profile.

As a second example, we add a square optical lattice to the same
harmonic confinement,
\begin{equation}
\label{eq:2d-optical-lattice}
V_{\rm ol}(x,y)
=
V_{\rm ah}(x,y)
+
V_0
\left[
\sin^2(kx)+\sin^2(ky)
\right],
\end{equation}
with $V_0=5$ and $k=\pi/2$. The additional periodic structure creates
multiple local wells and leads to the more structured density profile
shown in the right panel of
Fig.~\ref{fig:two-dimensional-examples}. The same potential-proximal
formulation applies directly to multi-well trapping potentials.

\begin{figure}[H]
\centering
\includegraphics[width=0.48\textwidth]
{figs/ground_state_2d_anisotropic_harmonic.eps}
\hfill
\includegraphics[width=0.48\textwidth]
{figs/ground_state_2d_optical_lattice.eps}
\caption{
Two-dimensional ground-state densities computed on $[-16,16]^2$
with $N_x=N_y=512$, with the central region shown.
The left panel corresponds to the anisotropic harmonic trap
\eqref{eq:2d-anisotropic-harmonic}, while the right panel includes
the square optical lattice \eqref{eq:2d-optical-lattice}.
Both computations use $\eps=10^{-2}$, $\sig=10^{-4}$, and
$\beta=\delta=10$.
}
\label{fig:two-dimensional-examples}
\end{figure}


% ============================================================
\section{Conclusion}
\label{sec:summary}
% ============================================================

%We have developed a positive-conservative Fourier optimization method
%for computing ground states of Bose--Einstein condensates with
%higher-order interactions in the density formulation. The singular
%kinetic term and the behavior of the potential term near zero density
%are treated by two compatible convex regularizations. On the
%computational domain, the regularized energies $\Gamma$-converge to the
%original density energy as the regularization parameters vanish.
%
%The resulting discretization preserves mass and nodal positivity.
%A potential-proximal FISTA iteration provides the main minimization,
%followed by a Newton refinement for high-accuracy stationarity.
%Numerical experiments demonstrate spectral-type convergence for fixed
%smooth regularized problems and reveal the bias--resolution tradeoff
%associated with the potential regularization. A joint analysis of
%domain truncation, Fourier discretization, and regularization remains
%for future work.


We have developed a positive-conservative Fourier optimization method
for computing ground states of Bose--Einstein condensates with
higher-order interactions in the density formulation. The kinetic and
potential terms are regularized while preserving the convex structure
of the density energy. On the computational domain, the regularized
energies $\Gamma$-converge to the original density energy as the
regularization parameters vanish.

The numerical results show spectral-type convergence of the Fourier
discretization when the regularization parameters are fixed. 
They also show
that decreasing the potential regularization parameter reduces the
regularization error but requires a finer Fourier resolution to resolve
the sharper low-density transition. In the two-stage optimization,
FISTA provides the main energy reduction, while the subsequent Newton
refinement further reduces the stationarity residual to high accuracy.

A joint analysis of domain truncation, Fourier discretization, and
regularization remains for future work.


% ============================================================

\bigskip
\noindent{\bf Acknowledgements} X. Ruan acknowledges support from the National Natural Science Foundation of China under grant 12201436.




\appendix

%% ============================================================
%\section{KKT system for the positive mass-conservative projection}
%\label{app:projection-kkt}
%% ============================================================
%
%Consider the nodal projection
%\[
%\min_y
%\frac12\norm{y-z}_h^2
%\quad
%\text{subject to}
%\quad
%y\ge0,
%\qquad
%h\bm1^Ty=1.
%\]
%Its KKT conditions are
%\[
%y-z+\lambda\bm1-\mu=0,
%\qquad
%y\ge0,
%\qquad
%\mu\ge0,
%\qquad
%y_j\mu_j=0.
%\]
%It follows that
%\[
%y_j=\max\{z_j-\lambda,0\},
%\]
%where $\lambda$ is determined by the monotone mass equation
%\[
%h\sum_j\max\{z_j-\lambda,0\}=1.
%\]
%This gives the scalar characterization of the positive
%mass-conservative projection used in Section~\ref{sec:ps}.
%

% ============================================================
\section{Implementation of the Newton refinement}
\label{app:newton-implementation}
% ============================================================

For the Newton refinement, let
\[
a=r_\eps(\rho),\qquad
a'=r_\eps'(\rho),\qquad
a''=r_\eps''(\rho),\qquad
q=D_N\rho,\qquad
\dot q=D_N\eta.
\]
The products and quotients in the following formula are understood
componentwise. The Hessian action of the regularized kinetic term is
\begin{align}
\label{eq:hessian-action}
H_{\rm kin}(\rho)[\eta]
={}&
\frac14D_N^T
\left(
\frac{\dot q}{a}
-
\frac{q\,a'}{a^2}\eta
\right)
-
\frac18
\left[
2q\dot q\,\frac{a'}{a^2}
+
q^2
\left(
\frac{a''}{a^2}
-
\frac{2(a')^2}{a^3}
\right)\eta
\right].
\end{align}
The full Hessian action is
\begin{equation}
\label{eq:full-hessian-action}
H(\rho)[\eta]
=
H_{\rm kin}(\rho)[\eta]
+
Vp_\sig''(\rho)\eta
+
\beta\eta
+
\delta D_N^TD_N\eta.
\end{equation}

For the PCG preconditioner, let $D_{\rm FD}$ denote the periodic
forward-difference operator and define
\[
B_{\rm FD}
=
D_{\rm FD}
-
\operatorname{diag}\left(\frac{q\,a'}{a}\right),
\qquad
W
=
\operatorname{diag}\left(\frac{1}{4a}\right),
\qquad
R_{\rm kin}
=
\operatorname{diag}\left(
-\frac{q^2a''}{8a^2}
\right).
\]
We use the sparse matrix
\begin{equation}
\label{eq:fd-preconditioner}
P_{\rm FD}
=
B_{\rm FD}^TWB_{\rm FD}
+
R_{\rm kin}
+
\delta D_{\rm FD}^TD_{\rm FD}
+
\operatorname{diag}\bigl(\beta+Vp_\sig''(\rho)\bigr)
\end{equation}
as the PCG preconditioner. In the strongly convex regime,
$P_{\rm FD}$ is symmetric positive definite.

For the shifted denominator used in the Newton refinement,
$r_\eps(s)=s+\eps$, we have $a'=1$ and $a''=0$.

\begin{thebibliography}{99}

\bibitem{AltmannHenningPeterseim2021}
R.~Altmann, P.~Henning, and D.~Peterseim,
\newblock The J-method for the Gross--Pitaevskii eigenvalue problem,
\newblock \emph{Numer. Math.}, 148 (2021), pp.~575--610.

\bibitem{Anderson1995}
M.~H. Anderson, J.~R. Ensher, M.~R. Matthews, C.~E. Wieman,
and E.~A. Cornell,
\newblock Observation of Bose--Einstein condensation in a dilute
atomic vapor,
\newblock \emph{Science}, 269 (1995), pp.~198--201.

\bibitem{AntoineDuboscq2014}
X.~Antoine and R.~Duboscq,
\newblock Robust and efficient preconditioned Krylov spectral solvers
for computing the ground states of fast rotating and strongly
interacting Bose--Einstein condensates,
\newblock \emph{J. Comput. Phys.}, 258 (2014), pp.~509--523.

\bibitem{AntoineLevittTang2017}
X.~Antoine, A.~Levitt, and Q.~Tang,
\newblock Efficient spectral computation of the stationary states of
rotating Bose--Einstein condensates by preconditioned nonlinear
conjugate gradient methods,
\newblock \emph{J. Comput. Phys.}, 343 (2017), pp.~92--109.

\bibitem{BaoCaiReview}
W.~Bao and Y.~Cai,
\newblock Mathematical theory and numerical methods for Bose--Einstein condensation,
\newblock \emph{Kinet. Relat. Models}, 6 (2013), pp.~1--135.

\bibitem{BaoCaiRuan2019}
W.~Bao, Y.~Cai, and X.~Ruan,
\newblock Ground states of Bose--Einstein condensates with higher order interaction,
\newblock \emph{Physica D}, 386--387 (2019), pp.~38--48.

\bibitem{BaoChernLim2006}
W.~Bao, I.-L.~Chern, and F.~Y.~Lim,
\newblock Efficient and spectrally accurate numerical methods for
computing ground and first excited states in Bose--Einstein
condensates,
\newblock \emph{J. Comput. Phys.}, 219 (2006), pp.~836--854.

\bibitem{BaoDu2004}
W.~Bao and Q.~Du,
\newblock Computing the ground state solution of Bose--Einstein condensates by a
normalized gradient flow,
\newblock \emph{SIAM J. Sci. Comput.}, 25 (2004), pp.~1674--1697.

\bibitem{BaoRuan2019}
W.~Bao and X.~Ruan,
\newblock Computing ground states of Bose--Einstein condensates with higher order
interaction via a regularized density function formulation,
\newblock \emph{SIAM J. Sci. Comput.}, 41 (2019), pp.~B1284--B1309.

\bibitem{BaoTang2003}
W.~Bao and W.~Tang,
\newblock Ground-state solution of Bose--Einstein condensate by
directly minimizing the energy functional,
\newblock \emph{J. Comput. Phys.}, 187 (2003), pp.~230--254.

\bibitem{BeckTeboulle2009}
A.~Beck and M.~Teboulle,
\newblock A fast iterative shrinkage-thresholding algorithm for linear inverse problems,
\newblock \emph{SIAM J. Imaging Sci.}, 2 (2009), pp.~183--202.

\bibitem{Braides2002}
A.~Braides,
\newblock \emph{Gamma-Convergence for Beginners},
\newblock Oxford University Press, Oxford, 2002.


\bibitem{lin2024sinum}
Z. Cai, B. Lin and M. Lin,
\newblock A positive and moment-preserving {Fourier} spectral method,
\newblock \emph{SIAM J. Numer. Anal.}, 62 (2024), pp.~273--294.


\bibitem{ChambolleDossal2015}
A.~Chambolle and C.~Dossal,
\newblock On the convergence of the iterates of the ``fast iterative
shrinkage/thresholding algorithm'',
\newblock \emph{J. Optim. Theory Appl.}, 166 (2015), pp.~968--982.

\bibitem{ChenLuLuZhang2024}
Z.~Chen, J.~Lu, Y.~Lu, and X.~Zhang,
\newblock On the convergence of Sobolev gradient flow for the
Gross--Pitaevskii eigenvalue problem,
\newblock \emph{SIAM J. Numer. Anal.}, 62 (2024), pp.~667--691.

\bibitem{Chin2010}
C.~Chin, R.~Grimm, P.~Julienne, and E.~Tiesinga,
\newblock Feshbach resonances in ultracold gases,
\newblock \emph{Rev. Mod. Phys.}, 82 (2010), pp.~1225--1286.

\bibitem{CollinMassignanPethick2007}
A.~Collin, P.~Massignan, and C.~J. Pethick,
\newblock Energy-dependent effective interactions for dilute
many-body systems,
\newblock \emph{Phys. Rev. A}, 75 (2007), 013615.

\bibitem{DanailaKazemi2010}
I.~Danaila and P.~Kazemi,
\newblock A new Sobolev gradient method for direct minimization of the
Gross--Pitaevskii energy with rotation,
\newblock \emph{SIAM J. Sci. Comput.}, 32 (2010), pp.~2447--2467.

\bibitem{DanailaProtas2017}
I.~Danaila and B.~Protas,
\newblock Computation of ground states of the Gross--Pitaevskii
functional via Riemannian optimization,
\newblock \emph{SIAM J. Sci. Comput.}, 39 (2017), pp.~B1102--B1129.

\bibitem{FuWangGao2003}
H.~Fu, Y.~Wang, and B.~Gao,
\newblock Beyond the Fermi pseudopotential: A modified
Gross--Pitaevskii equation,
\newblock \emph{Phys. Rev. A}, 67 (2003), 053612.

\bibitem{HenningPeterseim2020}
P.~Henning and D.~Peterseim,
\newblock Sobolev gradient flow for the Gross--Pitaevskii eigenvalue
problem: Global convergence and computational efficiency,
\newblock \emph{SIAM J. Numer. Anal.}, 58 (2020), pp.~1744--1772.

\bibitem{Ruan2018}
X.~Ruan,
\newblock A normalized gradient flow method with
attractive--repulsive splitting for computing ground states of
Bose--Einstein condensates with higher-order interaction,
\newblock \emph{J. Comput. Phys.}, 367 (2018), pp.~374--390.

\bibitem{RuanCaiBao2016}
X.~Ruan, Y.~Cai, and W.~Bao,
\newblock Mean-field regime and Thomas--Fermi approximations of
trapped Bose--Einstein condensates with higher-order interactions in
one and two dimensions,
\newblock \emph{J. Phys. B}, 49 (2016), 125304.

\bibitem{TianCaiWuWen2020}
T.~Tian, Y.~Cai, X.~Wu, and Z.~Wen,
\newblock Ground states of spin-$F$ Bose--Einstein condensates,
\newblock \emph{SIAM J. Sci. Comput.}, 42 (2020), pp.~B983--B1013.

\bibitem{Ulbrich2011}
M.~Ulbrich,
\newblock \emph{Semismooth Newton Methods for Variational Inequalities and
Constrained Optimization Problems in Function Spaces},
\newblock SIAM, Philadelphia, 2011.

\bibitem{VekslerFishmanKetterle2014}
H.~Veksler, S.~Fishman, and W.~Ketterle,
\newblock Simple model for interactions and corrections to the
Gross--Pitaevskii equation,
\newblock \emph{Phys. Rev. A}, 90 (2014), 023620.

\bibitem{WuWenBao2017}
X.~Wu, Z.~Wen, and W.~Bao,
\newblock A regularized Newton method for computing ground states of
Bose--Einstein condensates,
\newblock \emph{J. Sci. Comput.}, 73 (2017), pp.~303--329.

\bibitem{ZinnerThogersen2009}
N.~T. Zinner and M.~Th{\o}gersen,
\newblock Stability of a Bose--Einstein condensate with higher-order
interactions near a Feshbach resonance,
\newblock \emph{Phys. Rev. A}, 80 (2009), 023607.

\end{thebibliography}

\end{document}
